% part: intuitionistic-logic
% chapter: propositions-as-types
% section: terms-to-proofs

\documentclass[../../../include/open-logic-section]{subfiles}

\begin{document}

\olfileid{int}{pty}{tp}

\olsection{从证明项恢复推导}

现在我们考虑另一个方向：将证明项~$N$ 翻译回 \Log{N2i} 中的一个推导。一般而言，这只能相对于某个上下文来进行，该上下文须为给定证明项中出现的每个自由变元~$x$ 包含配对 $x : !A$。但我们先从一个不含自由变元的例子开始，即项
\[
\lambd[\typeof{x}{!A \land
  !B}][\pair{\proj{2}{x}}{\proj{1}{x}}]
\]
它的形式为 $\lambd[\typeof{x}{!A \land !B}][N_1]$。由于
\Log{tN2i} 中唯一能生成此种形式项的规则是~\Intro{\lif}，可知由~$N$ 编码的推导必定以~\Intro{\lif} 结尾，即：
  \[
  \Axiom$x : !A\land !B \fCenter \pair{\proj{2}{x}}{\proj{1}{x}} : 
    !C$
  \RightLabel{\Intro\lif}
  \UnaryInf$\fCenter \lambd[\typeof{x}{!A \land
  !B}][\pair{\proj{2}{x}}{\proj{1}{x}}] : (!A \land !B) \lif !C$
  \DisplayProof
  \]
我们目前尚不知 $!C$ 为何，但确实知道前件为
$!A \land !B$，且前提的上下文必须包含
$x:!A \land !B$。

现在，前提对应的证明项已经确定：
$\pair{\proj{2}{x}}{\proj{1}{x}}$。该证明项编码的推导
必定以 \Intro{\land} 结尾。我们仍然不知道 $!C$ 是什么，但由于它是~\Intro{\land} 的结论，可知它必定是一个
合取式，即 $!C \ident (!C_1 \land !C_2)$。我们继续重构：
  \[
  \Axiom$x : !A\land !B \fCenter \proj{2}{x} : !C_1$
  \Axiom$x : !A\land !B \fCenter \proj{1}{x} : !C_2$
  \RightLabel{\Intro\land}
  \BinaryInf$x : !A\land !B \fCenter \pair{\proj{2}{x}}{\proj{1}{x}} : 
    !C_1 \land !C_2$
  \RightLabel{\Intro\lif}
  \UnaryInf$\fCenter \lambd[\typeof{x}{!A \land
  !B}][\pair{\proj{2}{x}}{\proj{1}{x}}] : (!A \land !B) \lif (!C_1 \land !C_2)$
  \DisplayProof
  \]
证明项 $\pi_2(x)$ 表明，左侧最顶部的相继式是由
\Elim{\land}[2] 得到的，即其前提具有 $!D
\land !C_1$ 的形式。在右侧，我们可以确定得出
$!C_2$ 的推理必定是 \Elim{\land}[1]，且其前提具有
$!C_2 \land !E$ 的形式。
  \[
  \Axiom$x : !A\land !B \fCenter x: !D\land !C_1$
  \RightLabel{\Elim\land[2]}
  \UnaryInf$x : !A\land !B \fCenter \proj{2}{x} : !C_1$
  \Axiom$x : !A\land !B \fCenter x: !C_2\land !E$
  \RightLabel{\Elim\land[1]}
  \UnaryInf$x : !A\land !B \fCenter \proj{1}{x} : !C_2$
  \RightLabel{\Intro\land}
  \BinaryInf$x : !A\land !B \fCenter \pair{\proj{2}{x}}{\proj{1}{x}} : 
    !C_1 \land !C_2$
  \RightLabel{\Intro\lif}
  \UnaryInf$\fCenter \lambd[\typeof{x}{!A \land
  !B}][\pair{\proj{2}{x}}{\proj{1}{x}}] : (!A \land !B) \lif (!C_1 \land !C_2)$
  \DisplayProof
  \]
由于最顶层相继式的证明项现在都是变元，我们
可知它们必定是公理。这就确定了 $!C_1$、$!C_2$、$!D$ 和
$!E$：必须有 $!D \ident !A$ 且 $!C_1 \ident !B$，否则左边相继式
就不是公理；同样必须有 $!C_2 \ident !A$ 且 $!E \ident !B$，否则
右边相继式也不是公理。至此，我们已完全恢复了该证明：
  \[
  \Axiom$x : !A\land !B \fCenter x: !A\land !B$
  \RightLabel{\Elim\land[2]}
  \UnaryInf$x : !A\land !B \fCenter \proj{2}{x} : !B$
  \Axiom$x : !A\land !B \fCenter x: !A\land !B$
  \RightLabel{\Elim\land[1]}
  \UnaryInf$x : !A\land !B \fCenter \proj{1}{x} : !A$
  \RightLabel{\Intro\land}
  \BinaryInf$x : !A\land !B \fCenter \pair{\proj{2}{x}}{\proj{1}{x}} : 
    !B \land !A$
  \RightLabel{\Intro\lif}
  \UnaryInf$\fCenter \lambd[\typeof{x}{!A \land
  !B}][\pair{\proj{2}{x}}{\proj{1}{x}}] : (!A \land !B) \lif (!B \land !A)$
  \DisplayProof
  \]


\end{document}